% !TeX root = part4.tex
\documentclass[../final.tex]{subfiles}
\begin{document}

% Конспект: 11.jpg, 22.jpg.
% Условие: /Задачи/ТиСФ2_семинар_4.pdf на https://disk.yandex.ru/d/o4IQOWgLr_1y-g
\practice{4}{25.09.2026}{Полупроводники}

\rtheory{Электроны и дырки в собственном полупроводнике}

Полагаем $k_B=1$. Нуль энергии выберем на вершине валентной зоны;
дно зоны проводимости находится на высоте $\Delta$.
При $T=0$ в собственном полупроводнике валентная зона заполнена,
а зона проводимости пуста. Вблизи краёв зон используем параболические
законы дисперсии:
\[
\varepsilon_c(p)=\Delta+\frac{p^2}{2m_e^*},\qquad
    \varepsilon_v(p)=-\frac{p^2}{2m_h^*}.
\]
Здесь $m_e^*$ и $m_h^*$ --- эффективные массы электронов и дырок.
Рассматриваем по одной изотропной зоне каждого типа со спиновым
вырождением два. В реальном кристалле дополнительные долины учитываются
в эффективных плотностях состояний.

\begin{rbox}[left=16pt,right=16pt,top=12pt,bottom=10pt]{[]}
\centering
\begin{tikzpicture}[>=Latex,every node/.style={text=rtext}]
    \draw[->] (0,-0.25)--(0,4.6) node[above] {$\varepsilon$};
    \foreach \y in {0.05,0.23,0.41,0.59,0.77,0.95}
        \draw[rconceptcolor] (0.3,\y)--(3.35,\y);
    \foreach \y in {3.05,3.23,3.41,3.59,3.77,3.95}
        \draw[rc3] (0.3,\y)--(3.35,\y);
    \node[left] at (0,0.95) {$0$};
    \node[left] at (0,3.05) {$\Delta$};
    \draw[<->] (0.65,1.1)--(0.65,2.9)
        node[midway,right] {$\Delta$};
    \node[anchor=west,align=left] at (3.6,3.5)
        {Зона проводимости\\\small пуста при $T=0$};
    \node[anchor=west,align=left] at (3.6,0.5)
        {Валентная зона\\\small заполнена при $T=0$};
    \node[anchor=west] at (3.6,2.05) {\small Запрещённая зона};
\end{tikzpicture}
\captionof{figure}{Зонная схема собственного полупроводника.
Тепловое возбуждение электрона оставляет в валентной зоне дырку.}
\end{rbox}

\subsection*{Число носителей в невырожденном приближении}

Вероятность заполнения состояния в зоне проводимости и вероятность
отсутствия электрона в валентной зоне равны
\[
f_e(p)=\frac{1}{e^{(\varepsilon_c(p)-\mu)/T}+1},\qquad
    f_h(p)=1-\frac{1}{e^{(\varepsilon_v(p)-\mu)/T}+1}.
\]
Учитывая два спиновых состояния, запишем
\[
N_e=\frac{2V}{(2\pi\hbar)^3}\int f_e(p)\,d^3p,\qquad
    N_h=\frac{2V}{(2\pi\hbar)^3}\int f_h(p)\,d^3p.
\]
При $\Delta-\mu\gg T$ и $\mu\gg T$ распределения имеют
больцмановский вид:
\[
f_e(p)\simeq e^{(\mu-\Delta)/T}e^{-p^2/(2m_e^*T)},\qquad
    f_h(p)\simeq e^{-\mu/T}e^{-p^2/(2m_h^*T)}.
\]
Поскольку
$\int e^{-p^2/(2m^*T)}\,d^3p=(2\pi m^*T)^{3/2}$, получаем
\begin{align*}
    N_e&=\mathcal N_c(T)e^{(\mu-\Delta)/T},&
    \mathcal N_c(T)&=2V\left(\frac{m_e^*T}{2\pi\hbar^2}\right)^{3/2},\\
    N_h&=\mathcal N_v(T)e^{-\mu/T},&
    \mathcal N_v(T)&=2V\left(\frac{m_h^*T}{2\pi\hbar^2}\right)^{3/2}.
\end{align*}
Обозначения $\mathcal N_c$ и $\mathcal N_v$ относятся к эффективному
числу состояний во всём объёме $V$, а не к числу занявших их носителей.

\subsection*{Химический потенциал собственного полупроводника}

В отсутствие примесей электрон и дырка возникают парой: $N_e=N_h$.
Условие электронейтральности даёт
\[
\mathcal N_c e^{(\mu-\Delta)/T}
    =\mathcal N_v e^{-\mu/T}.
\]
Отсюда
\[
\mu_i=\frac\Delta2+\frac T2\ln\frac{\mathcal N_v}{\mathcal N_c}
    =\frac\Delta2-\frac{3T}{4}\ln\frac{m_e^*}{m_h^*}.
\]
\[
N_i=N_e=N_h
    =\sqrt{\mathcal N_c\mathcal N_v}\,e^{-\Delta/(2T)}.
\]

При одинаковых эффективных массах $\mu_i$ находится в середине
запрещённой зоны. Формула описывает предел $T\to0$ из области $T>0$;
строго при $T=0$ условие заполнения допускает любое положение $\mu$
внутри щели.

\clearpage
\task[1]{Донорный полупроводник}

\condition
В кремний, элемент IV группы, добавлено небольшое количество атомов
фосфора из V группы. Четыре электрона примеси участвуют в валентных
связях, пятый слабо связан с донором. Донорный уровень $E_d$ расположен
ниже дна зоны проводимости $\Delta$; энергия ионизации
$\Delta-E_d\sim0{,}03$--$0{,}1\,\text{эВ}$. Уже при температурах порядка
$300\,\text{К}$ возможен переход этих электронов в зону проводимости.
Найти температурную зависимость химического потенциала и числа
электронов в зоне проводимости.

\begin{rbox}[left=16pt,right=16pt,top=12pt,bottom=10pt]{[]}
\centering
\begin{tikzpicture}[>=Latex,every node/.style={text=rtext}]
    \draw[->] (0,-0.15)--(0,4.5) node[above] {$\varepsilon$};
    \foreach \y in {0.1,0.3,0.5,0.7}
        \draw[rconceptcolor] (0.35,\y)--(3.5,\y);
    \foreach \y in {3.25,3.45,3.65,3.85}
        \draw[rc3] (0.35,\y)--(3.5,\y);
    \draw[rassumptioncolor,thick,dashed] (0.35,2.55)--(3.5,2.55);
    \node[left] at (0,0.7) {$0$};
    \node[left] at (0,2.55) {$E_d$};
    \node[left] at (0,3.25) {$\Delta$};
    \draw[->,thick,rresultcolor] (2.7,2.6)--(2.7,3.4);
    \fill[rtext] (2.7,2.55) circle (2pt);
    \node[anchor=west] at (3.8,3.55) {Зона проводимости};
    \node[anchor=west] at (3.8,2.55) {Донорный уровень};
    \node[anchor=west] at (3.8,0.4) {Валентная зона};
    \draw[<->] (1,2.6)--(1,3.2);
    \node[anchor=west] at (1.15,2.9) {\small $\Delta-E_d$};
\end{tikzpicture}
\captionof{figure}{Ионизация донора: связанный электрон переходит
с уровня $E_d$ в зону проводимости.}
\end{rbox}

\solution
Пусть $N_d$ --- полное число доноров, $N_d^+$ --- число ионизованных
доноров. Акцепторы отсутствуют. Считаем доноры независимыми,
а носители в зонах --- невырожденными. Введём обозначение
$\delta=\Delta-E_d>0$.

\subsection*{Статистическая сумма одного донора}

Донор может быть пустым (ионизованным) или содержать один связанный
электрон с одним из двух направлений спина. Двойное заполнение
исключаем. Большая статистическая сумма одного центра равна
\[
\mathcal Z_d=1+2e^{(\mu-E_d)/T},\qquad
    \Omega_d=-T\ln\mathcal Z_d.
\]
Среднее число связанных электронов на одном доноре находим
дифференцированием по химическому потенциалу:
\[
\overline n_d=-\frac{\partial\Omega_d}{\partial\mu}
    =\frac{2e^{(\mu-E_d)/T}}{1+2e^{(\mu-E_d)/T}}.
\]
Вероятность ионизации равна $1-\overline n_d$, поэтому
\[
N_d^+=\frac{N_d}{1+2e^{(\mu-E_d)/T}}.
\]
Множитель два в знаменателе обусловлен спиновым вырождением
нейтрального донора; его необходимо сохранять в низкотемпературном ответе.

\subsection*{Уравнение электронейтральности}

Положительный заряд создают ионизованные доноры и дырки, отрицательный
--- электроны в зоне проводимости. Следовательно,
\[
N_e=N_d^++N_h.
\]
Подставляя найденные распределения, получаем основное уравнение
для химического потенциала:
\[
\mathcal N_c e^{(\mu-\Delta)/T}
    =\frac{N_d}{1+2e^{(\mu-E_d)/T}}+\mathcal N_v e^{-\mu/T}.
\]
После нахождения $\mu(T)$ число электронов определяется формулой
$N_e(T)=\mathcal N_c(T)e^{(\mu(T)-\Delta)/T}$.

Уравнение определяет переходы между режимами в рамках принятого
приближения. Рассмотрим его предельные случаи.

\subsection*{Низкие температуры: вымораживание носителей}

При $T\ll\delta$ большинство доноров нейтрально:
$2e^{(\mu-E_d)/T}\gg1$. Дырок мало, $N_h\ll N_e$, и
\[
\mathcal N_c e^{(\mu-\Delta)/T}
    \simeq\frac{N_d}{2}e^{(E_d-\mu)/T}.
\]
Логарифмируя, получаем
\[
\mu(T)\simeq\frac{\Delta+E_d}{2}
    -\frac T2\ln\frac{2\mathcal N_c(T)}{N_d}.
\]
Число свободных электронов равно
\[
N_e(T)\simeq\sqrt{\frac{N_d\mathcal N_c(T)}{2}}
    \,e^{-\delta/(2T)}.
\]
Таким образом, $N_e\propto T^{3/4}e^{-\delta/(2T)}$, а
$\mu\to(\Delta+E_d)/2$ при $T\to0$. Условие малой доли ионизованных
доноров можно записать как
$\mathcal N_c e^{-\delta/T}\ll2N_d$.

\subsection*{Переход к полной ионизации доноров}

Пока дырками можно пренебречь, $N_e=N_d^+$. Исключая $\mu$ с помощью
$e^{(\mu-E_d)/T}=(N_e/\mathcal N_c)e^{\delta/T}$, получаем
\[
\frac{2e^{\delta/T}}{\mathcal N_c}N_e^2+N_e-N_d=0.
\]
Положительный корень удобно представить в виде
\[
N_e(T)=\frac{2N_d}
    {1+\sqrt{1+8(N_d/\mathcal N_c(T))e^{\delta/T}}}.
\]
Соответствующий химический потенциал равен
\[
\mu(T)=\Delta+T\ln\frac{N_e(T)}{\mathcal N_c(T)}.
\]
Эта формула непрерывно связывает вымораживание с областью истощения
примеси. Когда
$2(N_d/\mathcal N_c)e^{\delta/T}\ll1$ и собственных носителей ещё мало,
каждый донор отдаёт практически один электрон:
\[
N_e\simeq N_d,\qquad
    \mu\simeq\Delta+T\ln\frac{N_d}{\mathcal N_c(T)}.
\]
\rnote{Граница полной ионизации зависит и от числа доноров, и от
плотности состояний; одного сравнения $T$ с $\delta$ недостаточно.}

\subsection*{Высокие температуры: вклад собственных носителей}

При полной ионизации доноров, но уже заметном числе дырок,
условие нейтральности имеет вид
\[
\mathcal N_c e^{(\mu-\Delta)/T}=N_d+\mathcal N_v e^{-\mu/T}.
\]
Положим $x=e^{\mu/T}$. Тогда
\[
\mathcal N_c e^{-\Delta/T}x^2-N_dx-\mathcal N_v=0,
\]
откуда
\[
e^{\mu/T}=\frac{N_d+
    \sqrt{N_d^2+4\mathcal N_c\mathcal N_v e^{-\Delta/T}}}
    {2\mathcal N_c e^{-\Delta/T}}.
\]
Введённое выше собственное число носителей удовлетворяет
$N_i^2=\mathcal N_c\mathcal N_v e^{-\Delta/T}=N_eN_h$.
Поэтому результат можно записать компактно:
\[
N_e=\frac{N_d+\sqrt{N_d^2+4N_i^2}}{2},\qquad
    N_h=\frac{\sqrt{N_d^2+4N_i^2}-N_d}{2}.
\]
\[
\mu=\Delta+T\ln\frac{N_d+\sqrt{N_d^2+4N_i^2}}
    {2\mathcal N_c}.
\]

При $N_i\ll N_d$ восстанавливается $N_e\simeq N_d$.
В собственной области $N_i\gg N_d$ получаем
\[
N_e\simeq N_h\simeq N_i,\qquad
    \mu\simeq\mu_i=\frac\Delta2-\frac{3T}{4}\ln\frac{m_e^*}{m_h^*}.
\]
Здесь <<высокая температура>> означает выход в собственную область;
невырожденное приближение всё ещё требует $\mu\gg T$ и $\Delta-\mu\gg T$.

\clearpage
\answer
\rconcept{Вымораживание:} свободных электронов экспоненциально мало,
\[
N_e\simeq\sqrt{\frac{N_d\mathcal N_c}{2}}\,e^{-\delta/(2T)},\qquad
    \mu\simeq\frac{\Delta+E_d}{2}-\frac T2\ln\frac{2\mathcal N_c}{N_d}.
\]
\rconcept{Истощение примеси:} доноры ионизованы, дырок мало,
\[
N_e\simeq N_d,\qquad \mu\simeq\Delta+T\ln\frac{N_d}{\mathcal N_c}.
\]
\rconcept{Собственная проводимость:} тепловых пар намного больше, чем доноров,
\[
N_e\simeq\sqrt{\mathcal N_c\mathcal N_v}\,e^{-\Delta/(2T)},\qquad
    \mu\simeq\frac\Delta2-\frac{3T}{4}\ln\frac{m_e^*}{m_h^*}.
\]

Все формулы для концентраций получаются делением чисел частиц на
$V$. Для температуры в кельвинах в формулах следует заменить $T$ на $k_BT$.

\end{document}
